\begin{frame}{확률 없이: 사람(메일)을 세는 방식}
\begin{center}\begin{varwidth}{0.92\textwidth}\usebeamercolor[fg]{normal text}
  
  \begin{tabular}{@{}lcc@{}}
    \toprule
     & 스팸 (300통) & 정상 (700통) \\
    \midrule
    `무료' \textbf{들어감} & $300\times0.6=$ \textbf{180} & $700\times0.05=$ \textbf{35} \\
    `무료' 없음 & 120 & 665 \\
    \bottomrule
  \end{tabular}
  \vskip0.8em
  \begin{itemize}
    
    \item $P(\text{spam}|\text{무료}) = 180/215 \approx 0.837$
      --- \textbf{빈도를 센 것만으로} 같은 답.
    \item \textbf{확률 트리: 사후확률 계산}
    
  \end{itemize}
\end{varwidth}\end{center}
\note{\begin{itemize}
\item 1000통의 메일(스팸 300 · 정상 700)을 세어 보자:
  \small
\item `무료'가 든 메일이라는 \textbf{조건}으로 세상을 좁히면
      총 215통(180+35), 그중 스팸 180통.
\item \kb{확률 트리}: \textbf{세로 합(215)이 조건부 세계의 전체
      크기}, 그 안에서 원하는 칸(180)의 비율이 곧 사후확률.
\item 의학과 보험, 사기 탐지 등 모든 베이즈 문제의
      \textbf{검증 도구}.
\end{itemize}}
\end{frame}
